% !TEX root = ../SANER2027-main-labeling.tex
% Related Work: each paragraph describes a line of prior work, then contrasts it with this study.
% Facts about the cited papers (quotes, pages): docs/reframe/RW_DOSSIER.md

\section{Related Work}
\label{sec:02_related}

\myparagraph{Supervised Issue Report Classification} Early \irc\ studies trained classic machine-learning classifiers on issue text~\cite{antoniol2008bug, kallis2019ticket, fan2017road}. Later work fine-tuned transformer encoders~\cite{colavito2022issue, bharadwaj2022github, trautsch2022predicting, izadi2022catiss, izadi2022predicting}. For instance, Izadi et al.~\cite{izadi2022catiss} and Trautsch et al.~\cite{trautsch2022predicting} fine-tuned RoBERTa~\cite{liu2019roberta} and seBERT~\cite{von2022validity}, respectively, to classify issues as bugs, features, or questions, and Siddiq and Santos trained a BERT classifier on more than 800{,}000 labeled issues~\cite{siddiq2022bert}. Colavito et al. fine-tuned a sentence encoder with SetFit on about 200 labeled issues~\cite{colavito2023few, colavito2024large}. All of these methods train a dedicated classifier on the labeled issues. This study instead holds the LLM fixed and varies only how the labeled issues are used, as examples in the prompt or as training data for LoRA, so its baseline is LoRA fine-tuning of the same models rather than a different classifier. Our sentence encoder is not trained and serves only to retrieve the neighbors.

\myparagraph{LLMs for Issue Report Classification} LLMs can classify issues without task-specific training~\cite{heo2025study, colavito2024leveraging, colavito2025benchmarking}, although their zero-shot performance varies across datasets~\cite{colavito2025benchmarking}. Colavito et al.~\cite{colavito2025benchmarking, colavito2024leveraging} added one or two randomly chosen labeled examples per class to the prompt, and in their benchmark these examples did not improve over zero-shot prompting for most models~\cite{colavito2025benchmarking}. The best reported results come from GPT models fine-tuned on labeled issues through OpenAI's API~\cite{aracena2024applyinglargelanguagemodels, aracena2025applying, heo2025study}. Open models such as Llama-3.1-8B and DeepSeek-R1-Distill-Llama-8B were fine-tuned with LoRA~\cite{hu2022lora} and scored lower~\cite{heo2025study, aracena2025applying}. Heo and Lee trained each classifier either on one project's issues or on the issues pooled from eleven projects~\cite{heo2025study}. We adopt their benchmark and both settings. These studies report cost as training time, inference time, or API prices~\cite{heo2025study, aracena2025applying, colavito2025benchmarking}, not as GPU memory. For future work, Heo and Lee suggest few-shot prompting, and Aracena et al. suggest retrieving relevant past issues to reduce the need for extensive fine-tuning~\cite{heo2025study, aracena2025applying}. This study evaluates both suggestions together. We show the LLM the labeled issues most similar to the query, up to 15 of them, and compare this with LoRA fine-tuning of the same four open models on the same labeled issues, in both of Heo and Lee's settings, including peak GPU memory. \rag\ outperforms zero-shot prompting at every model size and every $k$ (\Cref{sec:rq1}). To our knowledge, retrieval-augmented few-shot prompting has not been systematically evaluated for \irc\ or compared with LoRA fine-tuning before.

\myparagraph{Retrieved In-Context Examples} In-context learning lets an LLM infer a task from labeled examples in its prompt~\cite{brown2020language}. Choosing the examples most similar to the query outperforms random choice for GPT-3 on NLP benchmarks~\cite{liu2022makes}, and Milios et al. retrieve examples with a frozen sentence encoder to classify intents and emotions with up to 150 labels~\cite{milios2023context}. Yu et al., by contrast, train both the retriever and a RoBERTa classifier~\cite{yu2023retrieval}. In software engineering, LLM-Cure assigns app reviews to features with five fixed examples in its prompt~\cite{assi2026llm}. Closest to our work, Dinç and Tüzün show LLMs the five most similar labeled reports, retrieved with the same sentence encoder and FAISS library as ours, to decide whether Firefox bug reports are valid. In their study, fine-tuned encoders score higher than the prompted LLMs, and the authors name a sweep over $k$ and filtering the neighbors as future work~\cite{dincc2025judge}. For masked language models of under one billion parameters, Logan et al. found fine-tuning more accurate than in-context learning and more robust to the prompt wording~\cite{logan2021cuttingpromptsparameterssimple}. This study brings retrieved examples to three-label \irc\ across eleven projects. It sweeps $k$ up to where a vote over the neighbors' labels peaks, separates what retrieval and the LLM contribute with the \knn\ and zero-shot references (\Cref{sec:rq1}), and compares with LoRA fine-tuning of the same 3B to 32B instruction-tuned LLMs. In our study, \rag\ trails fine-tuning on the pooled issues by 2.8 points; with the example filter, the gap shrinks to 1.0 point and is not statistically significant (\Cref{sec:rq3}).

\myparagraph{Label Bias in the Examples} Ma et al. observe that similarity-based selection tends to pick examples whose labels match the query's, and that LLMs then lean toward the labels that dominate the context~\cite{ma2023fairnessguidedfewshotpromptinglarge}. Milios et al. report that the neighbors retrieved for the neutral emotion rarely carry that label, and suspect that retrieval limits performance on it~\cite{milios2023context}. We see a similar pattern for questions. On our validation split, true questions retrieve nearly as many bug-labeled as question-labeled neighbors (\Cref{sec:approach-bragtag}), and the models label many questions as bugs even without examples (\Cref{sec:rq1}). We hypothesize that the bug-labeled examples reinforce this error. Ma et al. reduce the bias by searching, with extra LLM calls, for one low-bias set of examples, scored by the entropy of the LLM's label distribution on a content-free input~\cite{ma2023fairnessguidedfewshotpromptinglarge}. Dinç and Tüzün add a judge LLM that sees the votes of several classifiers~\cite{dincc2025judge}. \Frag\ instead keeps per-query retrieval and applies a count rule to the neighbors' labels, removing the bug-labeled examples for suspected questions without an extra LLM call or access to output probabilities. It raises macro~$F_1$ over \rag\ by 1.5--2.4 points at the cost of bug recall (\Cref{sec:rq2}).

% Prior-work context for the question-to-bug confusion (ESEM reviewer C).
% Evidence, quotes and page numbers: docs/research/QUESTION_BUG_PRIOR_WORK.md
\myparagraph{Question-to-Bug Misclassification} In \irc\ studies that include both bug and question labels, question is often the hardest class to identify, and misclassified questions are most often labeled as bugs. For example, almost all methods in the NLBSE'24 tool competition struggled with question~\cite{kallis2024nlbse}, and on the same dataset, the two best-performing zero-shot LLMs labeled 46--54\% of questions as bugs~\cite{colavito2025benchmarking}. Fine-tuned LLMs also score lowest on question in aggregate~\cite{aracena2025applying, heo2025study}. Our results show the same pattern: every method we evaluate, fine-tuning included, labels many questions as bugs (\Cref{sec:rq1,sec:rq3}). As remedies, Aracena et al.~\cite{aracena2025applying} suggest retrieval augmentation, particularly for ambiguous labels such as question, and a first classification stage that separates questions from bug and feature reports. In our study, retrieved examples lower the share of questions labeled as bugs from 46--59\% under zero-shot prompting to 26--36\%, and \frag\ lowers it further to 19--27\% without an extra classification stage, but neither removes the confusion.

Prior work offers several explanations for this confusion. Question is often the smallest class in the training data~\cite{izadi2022catiss, bharadwaj2022github}, but our benchmark is balanced and the confusion remains, so imbalance alone does not explain it. Some questions read like bug reports because they describe unexpected behavior, contain error messages~\cite{colavito2024impact}, or lack explicit interrogative markers~\cite{aracena2025applying}. Our failure analysis (\Cref{sec:disc-failures}) finds a related cause: questions filed under a bug-report template lead the model to predict bug regardless of their content. Finally, a label can depend on information outside the report, such as a maintainer's later diagnosis~\cite{colavito2024impact}. Likewise, the label-noise cases in our failure analysis have labels that do not match their content and cannot be inferred from the report text.
